% Generated draft: requires mathematical and source review.
\documentclass[11pt]{article}
\usepackage[margin=25mm]{geometry}
\usepackage{fontspec}
\setmainfont{Latin Modern Roman}
\usepackage{amsmath,amssymb,mathtools,unicode-math}
\setmathfont{Latin Modern Math}
\usepackage{enumitem,needspace,xurl,hyperref}
\setlength{\parindent}{0pt}
\setlength{\parskip}{5pt}
\setlength{\emergencystretch}{3em}
\newcommand{\N}{\mathbb{N}}
\newcommand{\Z}{\mathbb{Z}}
\newcommand{\Q}{\mathbb{Q}}
\newcommand{\R}{\mathbb{R}}
\newcommand{\C}{\mathbb{C}}
\newcommand{\F}{\mathbb{F}}
\newcommand{\A}{\mathbb{A}}
\newcommand{\PP}{\mathbb{P}}
\newcommand{\E}{\mathbb{E}}
\newcommand{\eps}{\varepsilon}
\newcommand{\dd}{\,\mathrm{d}}
\begin{document}
\section*{UnsolvedMath ID: 801174}
\textbf{Source problem code:} catalogue entry
\subsection{Definitions and mathematical statement}
The supplied excerpt does not specify the ambient dimension or the dilation and translation conventions. The following statement uses the conventional one-dimensional dyadic interpretation of an orthonormal wavelet, with complex-valued functions; these conventions require confirmation against the full source.

Let \(\N_0=\{0,1,2,\ldots\}\). Define the Schwartz class by
\[
\mathcal S(\R)=\left\{f\in C^\infty(\R;\C):\sup_{x\in\R}(1+|x|)^m|f^{(n)}(x)|<\infty\text{ for every }m,n\in\N_0\right\}.
\]
For \(\psi\in\mathcal S(\R)\), use the Fourier transform convention
\[
\widehat\psi(\xi)=\int_{\R}\psi(x)e^{-2\pi i x\xi}\,\dd x.
\]
Call \(\psi\) an orthonormal dyadic wavelet if
\[
\left\{\psi_{j,k}:j,k\in\Z\right\},\qquad
\psi_{j,k}(x)=2^{j/2}\psi(2^j x-k),
\]
is an orthonormal basis of \(L^2(\R;\C)\).

\textbf{Question (Bownik and Rzeszotnik, Problem 13, 2019).} Is it true that every orthonormal dyadic wavelet \(\psi\in\mathcal S(\R)\) has compactly supported Fourier transform? Equivalently, does every such \(\psi\) satisfy
\[
\exists B>0\quad\forall\xi\in\R,\qquad |\xi|>B\Longrightarrow\widehat\psi(\xi)=0?
\]
The source motivates the question by stating that the only known construction of Schwartz-class wavelets is a Meyer wavelet construction, which is band-limited. This is the supplied source's historical qualification, not an independently verified assertion about all constructions or the current literature. The registry labels the problem open; that label is supplied metadata, not an independently verified mathematical assessment.
\subsection{Short English statement}
Under the explicitly stated one-dimensional dyadic convention, must every orthonormal wavelet that is smooth and rapidly decreasing have a Fourier transform that vanishes outside a bounded interval? The excerpt leaves the wavelet conventions implicit, so this interpretation requires source confirmation.
\subsection{Sources}
\begin{itemize}
\item Supplied local UnsolvedMath registry record, permanent ID \(801174\), source problem code \texttt{catalogue entry}, title \emph{Open Problems in Wavelet Theory --- Problem 13}. The record attributes the problem to Bownik and Rzeszotnik, with proposed year 2019, and labels it open.
\item Bownik and Rzeszotnik, \emph{Open Problems in Wavelet Theory} (2019), Problem 13. Supplied source URL: \url{https://arxiv.org/abs/1905.08230}. The record reports access on 2026-10-16 and extraction by \texttt{arxiv\_latex}. These are supplied provenance details; the source was not independently visited or verified for this draft.
\end{itemize}
\end{document}
